% Run GCT_KRON_20261120; GPT-6 Astra Pro; effort not exposed.
\section{An explicit unbounded gap between the determinant orbit and its closure}

\subsection*{A-028. Rectangular examples where symmetric Kronecker multiplicities lose copies}
For an integer \(h\geq2\), put
\[
 \rho_h=(h,h,h),\qquad \mu_h=(\underbrace{3,\ldots,3}_{h}),\qquad
 p_h=\#\{(a,b,c)\in\mathbb Z_{\geq0}^3:2a+3b+4c=h\}.
\]
Here the determinant has size \(n=h\), and the outer polynomial
 degree is \(d=3\). Then
\begin{align}
 g(\rho_h,\mu_h,\mu_h)&=p_h,                         \tag{A-028.1}\\
 \operatorname{sk}(\rho_h,\mu_h)&=
 \begin{cases}p_h&h\text{ even},\\0&h\text{ odd},\end{cases}
                                                               \tag{A-028.2}\\
 K_h(\rho_h)=a_{\rho_h}(3[h])&=
 \begin{cases}1&h\text{ even},\\0&h\text{ odd}.\end{cases}
                                                               \tag{A-028.3}
\end{align}
The symmetric Kronecker term means the multiplicity of the Specht
module \([\rho_h]\) in \(\operatorname{Sym}^2[\mu_h]\).
It is the determinant orbit multiplicity in the established
orbit-stabilizer formula recalled in B\"urgisser--Ikenmeyer--Panova,
\emph{No occurrence obstructions in geometric complexity theory},
Section 1, subsection ``Related work: Kronecker coefficients,
paragraph defining the GCT-coefficients (supplied source
\texttt{main.tex}, lines 415--428). Consequently, for every
even \(h\geq4\),
\[
 K_h(\rho_h)=1<\operatorname{sk}(\rho_h,\mu_h)
                       =g(\rho_h,\mu_h,\mu_h)=p_h.              \tag{A-028.4}
\]
These gaps are unbounded: \(p_{4t}\geq t+1\).
This is a concrete extraction from known invariant and plethysm
inputs, with the transpose action checked below; historical priority
is not asserted. It exhibits why an orbit multiplicity cannot be
substituted for the actual closure multiplicity.

\paragraph{The transpose action on the known invariant generators.}
Use the algebra \(\mathcal A\) and generators \(H,Q,\widetilde S\)
from A-026. Identify all three three-dimensional vector spaces with
\(\mathbb C^3\), with fixed ordered bases. Regard a tensor as a
triple of \(3\)-by-\(3\) matrices \((A_1,A_2,A_3)\), and let
\(\tau\) simultaneously transpose the three matrices. This exchanges
the first two tensor factors and preserves \(\mathcal A\).
We claim
\[
              \tau H=H,\qquad\tau Q=-Q,\qquad
                     \tau\widetilde S=\widetilde S.             \tag{A-028.5}
\]
The degree-six and degree-nine components of the polynomial algebra
\(\mathcal A=\mathbb C[H,Q,\widetilde S]\) are respectively the
one-dimensional spaces spanned by \(H\) and by \(Q\). Thus \(\tau\)
acts on each by a scalar, necessarily \(\pm1\).

The primary source explicitly evaluates these two generators on
skew-symmetric matrix triples: Domokos--Drensky,
\emph{Defining relation for semi-invariants of three by three matrix
triples}, arXiv:1101.3178v1, Section 2, printed p.~5, the displayed
skew-symmetric triple immediately before the page break to p.~6.
Writing its three columns of skew coordinates as a matrix \(M\),
the displayed values are \(H=(\det M)^2\) and \(Q=(\det M)^3\).
Take \(M=I_3\). Both values are nonzero, while transposition sends
the entire triple to its negative. Homogeneity therefore multiplies
\(H\) by \((-1)^6=1\) and \(Q\) by \((-1)^9=-1\).
This determines the first two scalars in (A-028.5) on their entire
one-dimensional invariant spaces, not merely at this point.

For the last identity, each coefficient \(f_{i,j,k}\) of
\(\det(t_1A_1+t_2A_2+t_3A_3)\) is unchanged by simultaneous
transposition, since determinant is unchanged by transpose.
The definition of \(\widetilde S\) in the same source, Section 2,
printed p.~5, is a polynomial in these coefficients. It too is
unchanged. This completes the verification of (A-028.5).

\begin{proof}[Proof of all three multiplicity formulas]
By A-026, the degree-\(3h\) invariant space has basis
\[
                    H^a Q^b\widetilde S^c,
                    \qquad 2a+3b+4c=h.
\]
Every such \(b\) has the parity of \(h\), so (A-028.5) shows
that \(\tau\) acts on the entire space by \((-1)^h\).

In the three-factor Schur--Weyl decomposition (A-026.4), this
invariant space is the multiplicity space
\[
                  ([\rho_h]\otimes[\rho_h]\otimes[\rho_h])^{S_{3h}}.
                                                               \tag{A-028.6}
\]
The determinant lines in the three Schur factors are identified with
\(\mathbb C\) by the fixed ordered volume forms. Exchanging the
first two such lines acts trivially on their tensor product. These
are ordinary tensor spaces, with no graded or exterior sign in the
interchange. Thus the induced action on (A-028.6) is precisely the
flip of its first two Specht factors. Its fixed space has dimension
\(p_h\) for even \(h\), and zero for odd \(h\).

Conjugating a partition twists its Specht module by the sign
representation. Since \(\mu_h=\rho_h'\), one has
\([\mu_h]\cong[\rho_h]\otimes\operatorname{sgn}\).
The natural cancellation of the two sign factors gives an
\(S_{3h}\)-equivariant isomorphism
\[
 [\mu_h]\otimes[\mu_h]\ \cong
 [\rho_h]\otimes[\rho_h]
\]
which commutes with the flip. In particular their symmetric squares
are isomorphic as \(S_{3h}\)-modules. Taking the \([\rho_h]\)
multiplicity, using characteristic-zero self-duality of Specht modules,
proves (A-028.2). Taking the full tensor multiplicity proves
(A-028.1), also directly the usual double-conjugation invariance
of Kronecker coefficients. The precise classical sign-twist and
Schur--Weyl inputs are the character involution in Macdonald,
Chapter I, Section 7, equations (7.3)--(7.6), pp.~113--114, and
Appendix A(5.3)--(5.6), pp.~155--156.

For the ambient coefficient, apply F-011 with \(d=3,n=k=h\).
Its partition \((d(n-k)+k,k,k)\) is exactly \(\rho_h\),
and its value is one for even \(h\), zero for odd \(h\).
Since \(\rho_h\) has three rows and \(h\geq2\), O-007 gives
\(K_h(\rho_h)=a_{\rho_h}(3[h])\). This proves (A-028.3).
O-007 uses Beauville, \emph{Determinantal hypersurfaces},
Proposition 3.1(a), printed p.~46; all of that input's hypotheses
and the degree-two endpoint were checked there.

For even \(h\geq4\), the two solutions
\((a,b,c)=(h/2,0,0)\) and \(((h-4)/2,0,1)\) are distinct,
so \(p_h\geq2\). For \(h=4t\), the \(t+1\) choices
\(b=0,c=0,\ldots,t,a=2t-2c\) prove the unbounded gap.
All uses of the invariant-ring theorem, ordinary/symmetric tensor
operations, and ambient-to-closure comparison are thereby separate.
\end{proof}

\paragraph{Consequences and their scope.}
This family is stronger than an example with \(K<g\) caused only
by the antisymmetric summand: at every even index all the \(p_h\)
copies already belong to the symmetric Kronecker term, while exactly
one extends to the determinant closure. The difference
\(p_h-1\) is a loss of orbit functions at the boundary.
It supplies no actual permanent obstruction: \(q\leq a=1=K\)
for even \(h\), and \(q=0\) for odd \(h\). Thus ordinary and
symmetric Kronecker coefficients, their positive monomial bases,
and their saturation behavior still do not compute the rank needed
for a general actual multiplicity obstruction.

\paragraph{Grothendieck dictionary.}
The tensor representations and the transposition automorphism are
maps of affine group-scheme representations over \(\mathbb Z\).
Their invariant spaces are kernels of coactions, as in SGA 3, I,
Proposition 4.7.2, pp.~34--35. The Schur decompositions and
symmetric-square summand are the characteristic-zero realization
of Schur operations on locally free modules from SGA 6, VI,
Theorem 3.3, pp.~392--393; the displayed calculation is at the
level of actual modules and their explicit involution, not only
of virtual classes. The degree gradings are SGA 3, I,
Proposition 4.7.3, pp.~35--36.

Let \(\Omega_h\) be the scheme-theoretic determinant orbit closure
and \(\mathcal O_h\) its dense orbit. Restriction is the injective
map of algebras
\(\mathbb C[\Omega_h]\hookrightarrow\mathbb C[\mathcal O_h]\).
The multiplicity subspaces of type \(\rho_h\) have respective
dimensions one and \(p_h\) at even \(h\). Therefore \(p_h-1\)
independent orbit copies do not extend as regular functions across
the closure. This boundary-extension condition is additional
geometric data, invisible in a positive class of \(K_0\).
The orbit-closure ideal and its graded quotients can be placed in
the Hilbert/Quot framework of FGA, Bourbaki expos\'e 221,
Theorem 3.1, p.~259, and Theorem 3.2, p.~260; representability
alone does not force these restriction maps to be surjective.
All multiplicity assertions here are in characteristic zero;
no arbitrary integral base-change assertion is made for invariants.
